\begingroup
\renewcommand{\arraystretch}{1.45}
\Needspace{40\baselineskip}
\InterpolationArticle{24}
Through the investigation presented from Article 20 onward, we have reached the point where, given $\mu$
values of the function $X$ of the variable angle $x$ that form a complete period, we can determine a function $(X')=$
\[
\begin{aligned}
&\gamma+\gamma'\cos x+\gamma''\cos2x+\mathinner{.\,.\,.}+\gamma^m\cos mx\\[5pt]
&\quad+\delta'\sin x+\delta''\sin2x+\mathinner{.\,.\,.}+\delta^m\sin mx
\end{aligned}
\]
that satisfies them all, with the number $m$ expressing the order of this function being either $=\frac12\mu-\frac12$ or $=\frac12\mu$, according as $\mu$ is odd or
even. From Articles 15 and 19 it follows that any other similar function satisfying the same values must be contained in the form $X'+Y\sin\frac12\mu(x-a)$,
where $Y$ denotes an undetermined function of the angle $x$ of the form $G$ or $F$ (Article 13), according as $\mu$ is
odd or even. In the former case, besides $X'$ there is no other similar function satisfying the same
values except one of higher order; hence $X'$ will be identical with $X$ if
it is known that $X$ does not exceed the $m^{\mathrm{th}}$ order. In the latter case, infinitely many
other similar functions of the $m^{\mathrm{th}}$ order do indeed satisfy the same values, but all agree in
being contained in the form $X'+h\sin m(x-a)$, with
$h$ denoting a fixed quantity. Consequently, if it is known that $X$ does not exceed the $m^{\mathrm{th}}$ order,
at least $X$ and $X'$ will not differ except in the coefficients of $\cos mx$ and $\sin mx$;
in the function $X'$, however, the equation
$\gamma^m\sin ma-\delta^m\cos ma=0$ will hold between these coefficients. Furthermore, it is clear that, if a function $X''$ similar in form to $X'$ and satisfying the same values is known from elsewhere, any function $X''+Y\sin\frac12\mu(x-a)$ satisfies them, and $Y$ can be determined so that this function does not exceed the $m^{\mathrm{th}}$ order. To this end we observe that, if a part of $Y$ is represented by $K\cos(\kappa x+k)$, with $K$
denoting a fixed coefficient, $k$ a fixed angle, and $\kappa$ a number that is either an integer or has fractional part
$\frac12$ (according as $\mu$ is even or odd)---a form to which the individual parts can be reduced by combining the cosine and sine of the same angle---the corresponding part of the product $Y\sin\frac12\mu(x-a)$ will be represented by
\SourceLeaf{311}
\[
\begin{aligned}
&\frac12K\sin((\kappa+\frac12\mu)x+k-\frac12\mu a)\\[5pt]&\qquad{}-\frac12K\sin((\kappa-\frac12\mu)x+k+\frac12\mu a)
\end{aligned}
\]
or, equivalently, by
\[
\begin{aligned}
&L\sin(\lambda x+l)-L\sin((\lambda-\mu)x+l+\mu a)
\end{aligned}
\]
\Needspace{10\baselineskip}
with $\lambda$ denoting an integer. Thus, in $X''$, any part $L\sin(\lambda x+l)$ may be replaced by $L\sin((\lambda-\mu)x+l+\mu a)$, and similarly this in turn by
$L\sin((\lambda-2\mu)x+l+2\mu a)$, then by $L\sin((\lambda-3\mu)x+l+3\mu a)$ etc. To extend a similar conclusion to cosines, it suffices to observe that cosines appear in place of sines
if only $l$ is replaced by $l+90^{0}$. It follows that, if in $X''$ a term
$L\cos\lambda x$ occurs, with $\lambda$ denoting an integer greater than $\frac12\mu$, which can therefore be put in the form
$\lambda=\nu\mu\pm\lambda'$ with $\nu$ an integer and $\lambda'$ no greater than $\frac12\mu$, this term may be replaced by
\[
\begin{aligned}
&L\cos(\pm\lambda'x+\nu\mu a)=L\cos\nu\mu a\cos\lambda'x\mp L\sin\nu\mu a\sin\lambda'x
\end{aligned}
\]
and similarly a term $L\sin\lambda x$ may be replaced by
\[
\begin{aligned}
&L\sin(\pm\lambda'x+\nu\mu a)=L\sin\nu\mu a\cos\lambda'x\pm L\cos\nu\mu a\sin\lambda'x
\end{aligned}
\]
In this way the function $X''$ will evidently be reduced to another whose order is certainly no greater
than the $\frac12\mu^{\mathrm{th}}$ order.

\Needspace{24\baselineskip}
By this operation, therefore, $X''$ will always become $X'$ itself whenever $\mu$
is odd; in the other case, where $\mu$ is even, it will certainly become a function that
can differ only in the coefficients of $\cos mx$ and $\sin mx$ from the function $X'$.
To derive the corresponding coefficients in $X'$ from those coefficients, the general method of Article 15 may be used. From this it follows without difficulty that, if
the last terms in that function are $K\cos mx+L\sin mx$, these must be replaced by the following in order to obtain the function $X'$:
\[
\begin{aligned}
&\quad(K\cos ma+L\sin ma)\cos ma\cos mx\\[5pt]
&+(K\cos ma+L\sin ma)\sin ma\sin mx
\end{aligned}
\]
These terms become $=K\cos mx+L\sin mx$ if between $K$ and $L$ the equation $K\sin ma=L\cos ma$
holds, as is readily confirmed by calculation; in that case, therefore, the function is already entirely identical with $X'$.
\endgroup
